% config.tex -- усі дані про вас і вашу роботу в одному місці.
%
% Це ЄДИНИЙ файл, який більшість студентів мають редагувати напряму.
% Він не залежить від факультету чи кафедри -- це просто текстові
% поля, тому шаблон однаково підходить для будь-якої кафедри КПІ:
% просто впишіть свої дані замість прикладів нижче.
%
% Ступінь (бакалавр/магістр) тут НЕ задається -- він вказаний у
% main.tex, в аргументі \documentclass[master]{kpithesis} або
% [bachelor]. Це свідомий поділ: тут -- дані про вас, там -- вибір
% "форми" документа.

% --- Заклад ---
\newcommand{\facultyName}{Повна назва факультету/навчально-наукового інституту}
\newcommand{\departmentName}{Повна назва кафедри}
\newcommand{\departmentHeadName}{Ім'я ПРІЗВИЩЕ}
% Завідувач(-ка) кафедри -- саме той підпис, що стоїть у грифі
% "ДО ЗАХИСТУ ДОПУЩЕНО" на титульній сторінці.

% --- Спеціальність і програма ---
\newcommand{\specialtyCode}{000}
\newcommand{\specialtyName}{Назва спеціальності}
\newcommand{\programmeName}{Назва освітньої програми}
% Для магістерської дисертації важливо, за якою саме програмою
% ви навчаєтесь -- освітньо-професійною чи освітньо-науковою (це
% впливає на формулювання на титульній сторінці). Впишіть одне з двох:
% "освітньо-професійною" або "освітньо-науковою". Для бакалаврської
% роботи це поле не використовується (там завжди ОПП) -- можете
% залишити як є.
\newcommand{\programmeType}{освітньо-професійною}

% --- Тема роботи ---
\newcommand{\thesisTitle}{Тема вашої роботи}
% Для бакалаврської роботи вкажіть одне з двох: "Дипломний проєкт" або
% "Дипломна робота" (перше -- якщо є графічна частина/креслення,
% друге -- якщо ні; для більшості неінженерних спеціальностей це
% "Дипломна робота"). Для магістра це поле не використовується --
% магістерська робота завжди називається "Магістерська дисертація".
\newcommand{\workType}{Дипломна робота}

% --- Автор(-ка) ---
\newcommand{\authorFullName}{Прізвище Ім'я По-батькові}
\newcommand{\authorCourse}{4}
\newcommand{\authorGroup}{XX-00}

% --- Науковий керівник ---
\newcommand{\supervisorFullName}{Прізвище Ім'я По-батькові}
\newcommand{\supervisorRegalia}{посада, науковий ступінь, вчене звання}
% Наприклад: "доцент кафедри ОТ, канд. техн. наук, доцент"

% --- Консультант (за наявності) ---
% Якщо консультанта немає -- просто залиште \consultantFullName
% порожнім (як зараз): титульна сторінка й бланк завдання самі
% пропустять цей рядок.
\newcommand{\consultantFullName}{}
\newcommand{\consultantRegalia}{}
\newcommand{\consultantChapter}{}
% \consultantChapter -- назва розділу, з якого консультували,
% наприклад "економічної частини" -- потрібна лише якщо консультант є.

% --- Рецензент ---
\newcommand{\reviewerFullName}{Прізвище Ім'я По-батькові}
\newcommand{\reviewerRegalia}{посада, науковий ступінь, вчене звання}

% --- Формальності ---
\newcommand{\udcCode}{000.000}
% УДК скаже Вам науковий керівник -- код потрібен
% лише для магістерської дисертації (для бакалаврської не друкується).
\newcommand{\thesisYear}{2026}
